% GENERATED FILE. Edit canonical lessons, then run npm run generate:books.
% canonical-source-hash: fnv1a64:01393a6765f628ef
% canonical-lessons: ZH-C408-zhanpai, ZH-C408-gongjiao, ZH-W408-jiao-cross, ZH-C408-toubi, ZH-C408-keche, ZH-C408-qidian

\chapter{By Bus}
\label{ch:zh-by-bus}
\hlchaptermodality{408}

\begin{chapteropening}
\textbf{By the end of this chapter:} \emph{I can find the bus stop, take public transport, pay with coins and say where a route starts.}

Complete the last lesson of chapter 408: I can find the bus stop, take public transport, pay with coins and say where a route starts.
\end{chapteropening}

\section[zhànpái]{zhànpái — bus stop sign}

\label{lesson:ZH-C408-zhanpai}



\begin{quote}
Before the new one: say the Mandarin for the end of the line, then the Mandarin for a D-class bullet train.
\end{quote}

\subsection*{You'll want to know: zhànpái}
\emph{zhànpái} — \textquotedblleft{}bus stop sign\textquotedblright{}.

\begin{quote}
\emph{wǒ zài zhànpái xiàbian děng nǐ} — I'll wait for you under the bus stop sign.
\end{quote}

Its characters come later; for now it is a word for the ear.

\begin{cousinweb}
\emph{zhànpái} is built from \zh{站} 'a stop' and \emph{pái} 'a board', both inherited from Middle Chinese, \emph{treamH} and \emph{bea} (Guangyun, in Baxter's transcription), which glosses the second as a signboard.
\end{cousinweb}

\begin{grammarlens}[title={What you've built}]
One more word for getting around.
\end{grammarlens}

\subsection*{Your turn}
\begin{itemize}
\raggedright
  \item \emph{Say it:} \emph{zhànpái}
  \item \emph{Say it:} \emph{zhànpái}, once more
  \item \emph{Say it:} \emph{wǒ zài zhànpái xiàbian děng nǐ} — I'll wait for you under the bus stop sign.
  \item \emph{Recall:} say the Mandarin for the end of the line, then the Mandarin for a D-class bullet train, then say \emph{zhànpái} again
\end{itemize}

\begin{tcolorbox}[breakable,colback=teal!4,colframe=teal!35!black,title={Before you move on}]
What does zhànpái mean? (\textquotedblleft{}Bus stop sign\textquotedblright{}.) Say it once more.
\end{tcolorbox}
\section[\texorpdfstring{\zh{公交} (gōngjiāo)}{公交 (gōngjiāo)}]{\zh{公交} (gōngjiāo) — public transport}

\label{lesson:ZH-C408-gongjiao}



\begin{quote}
Before the new one: say the Mandarin for a D-class bullet train, then the Mandarin for a bus stop sign.
\end{quote}

\subsection*{You'll want to know: \zh{公交}}
\textbf{\zh{公交}} — \emph{gōngjiāo} — \textquotedblleft{}public transport\textquotedblright{}.

\begin{quote}
\textbf{\zh{我每天坐公交去学校}} \emph{wǒ měitiān zuò gōngjiāo qù xuéxiào} — I take public transport to school every day.
\end{quote}

The new character, in pieces: \zh{交} — six strokes: a dot and a bar on top; under them a short stroke falling to the left and a short dot falling to the right; then a long stroke falling to the left, crossed by a long stroke falling to the right. The next lesson puts it in your hand.

\begin{cousinweb}
\textbf{\zh{公交}} (\emph{gōngjiāo}) is built from \zh{公} 'public' and \zh{交} 'to cross', both inherited from Middle Chinese, \emph{kuwng} and \emph{kaew} (Guangyun, in Baxter's transcription), which glosses the first 'shared' and the second 'to join'. It is a short form of a longer name for public transport; the city bus is \emph{gōngjiāochē}.
\end{cousinweb}

\begin{grammarlens}[title={What you've built}]
One more word for getting around.
\end{grammarlens}

\subsection*{Your turn}
\begin{itemize}
\raggedright
  \item \emph{Say it:} \emph{gōngjiāo}
  \item \emph{Say it:} \emph{gōngjiāo}, once more
  \item \emph{Say it:} \emph{wǒ měitiān zuò gōngjiāo qù xuéxiào} — I take public transport to school every day.
  \item \emph{Recall:} say the Mandarin for a D-class bullet train, then the Mandarin for a bus stop sign, then say \emph{gōngjiāo} again
\end{itemize}

\begin{tcolorbox}[breakable,colback=teal!4,colframe=teal!35!black,title={Before you move on}]
What does \zh{公交} mean? (\textquotedblleft{}Public transport\textquotedblright{}.) Say it once more.
\end{tcolorbox}
\section[\texorpdfstring{\zh{交} (jiāo)}{交 (jiāo)}]{\zh{交} — a lid, then two strokes crossing}

\label{lesson:ZH-W408-jiao-cross}



\begin{quote}
The word first, then the character inside it.

\begin{itemize}
\raggedright
  \item \emph{Recall:} say \emph{gōngjiāo} — \textbf{R1}, one lesson back, before it decays
\end{itemize}

\begin{quote}
\zh{文} \zh{公}
\end{quote}

\emph{Read it:} each one aloud, with nothing said first — \emph{wén}, \emph{gōng}
\end{quote}

\subsection*{Script — a lid, then two strokes crossing}
\hlblockfigure{\detokenize{figures/ZH-W408-jiao-cross-filmstrip.pdf}}{How \zh{交} is written, stroke by stroke}

\begin{quote}
\zh{交}
\end{quote}

\textbf{Two pieces, top to bottom.}

\textbf{The lid, two strokes.} A dot, then a long bar across, as on top of \zh{文}.

\textbf{The lower piece, four strokes.} A short stroke falling to the left and a short dot falling to the right, apart, like two feet. Then a long stroke falling to the left, and across it a long stroke falling to the right to the foot.

2 + 4 = six.

\subsection*{Your turn}
\begin{itemize}
\raggedright
  \item \emph{Write it:} \zh{交}, all six, with the model in front of you
  \item \emph{Cover:} write it again from memory, and count the pieces, not the strokes
\end{itemize}

\begin{tcolorbox}[breakable,colback=teal!4,colframe=teal!35!black,title={Before you move on}]
How many strokes in \zh{交}? (\textbf{Six}.) What sits on top? (\textbf{A dot and a bar, as on \zh{文}}.) What are the last two strokes? (\textbf{A long fall to the left, crossed by a long fall to the right}.)

Source: \href{https://raw.githubusercontent.com/chanind/hanzi-writer-data/68d10a4b21150cae5e1ebbd223eed289cf32d90c/data/\%E4\%BA\%A4.json}{Hanzi Writer Data \zh{交}, pinned stroke paths and medians}.
\end{tcolorbox}
\section[tóubì]{tóubì — to drop coins in the box}

\label{lesson:ZH-C408-toubi}



\begin{quote}
Before the new one: say the Mandarin for a bus stop sign, then the Mandarin for public transport.
\end{quote}

\subsection*{You'll want to know: tóubì}
\emph{tóubì} — \textquotedblleft{}to drop coins in the box\textquotedblright{}.

\begin{quote}
\emph{shàng chē yǐhòu qǐng tóubì} — Please put your coins in once you're on.
\end{quote}

Its characters come later; for now it is a word for the ear.

\begin{cousinweb}
\emph{tóubì} is built from \emph{tóu} 'to throw' and \emph{bì} 'money', both inherited from Middle Chinese, \emph{duw} and \emph{bjiejH} (Guangyun, in Baxter's transcription), which glosses the second as gifts of silk; the sense 'coin' is not given there.
\end{cousinweb}

\begin{grammarlens}[title={What you've built}]
One more word for getting around.
\end{grammarlens}

\subsection*{Your turn}
\begin{itemize}
\raggedright
  \item \emph{Say it:} \emph{tóubì}
  \item \emph{Say it:} \emph{tóubì}, once more
  \item \emph{Say it:} \emph{shàng chē yǐhòu qǐng tóubì} — Please put your coins in once you're on.
  \item \emph{Recall:} say the Mandarin for a bus stop sign, then the Mandarin for public transport, then say \emph{tóubì} again
\end{itemize}

\begin{tcolorbox}[breakable,colback=teal!4,colframe=teal!35!black,title={Before you move on}]
What does tóubì mean? (\textquotedblleft{}To drop coins in the box\textquotedblright{}.) Say it once more.
\end{tcolorbox}
\section[\texorpdfstring{\zh{客车} (kèchē)}{客车 (kèchē)}]{\zh{客车} (kèchē) — coach; long-distance bus}

\label{lesson:ZH-C408-keche}



\begin{quote}
Before the new one: say the Mandarin for public transport, then the Mandarin for to drop coins in the box.
\end{quote}

\subsection*{You'll want to know: \zh{客车}}
\textbf{\zh{客车}} — \emph{kèchē} — \textquotedblleft{}coach; long-distance bus\textquotedblright{}.

\begin{quote}
\textbf{\zh{我们坐客车去那儿}} \emph{wǒmen zuò kèchē qù nàr} — We're going there by coach.
\end{quote}

\begin{cousinweb}
\textbf{\zh{客车}} (\emph{kèchē}) is built from \zh{客} 'a guest' and \zh{车} 'vehicle', both inherited from Middle Chinese, \emph{khaek} and \emph{tsyhae} (Guangyun, in Baxter's transcription): a vehicle that carries people, not goods.
\end{cousinweb}

\begin{grammarlens}[title={What you've built}]
One more word for getting around.
\end{grammarlens}

\subsection*{Your turn}
\begin{itemize}
\raggedright
  \item \emph{Say it:} \emph{kèchē}
  \item \emph{Say it:} \emph{kèchē}, once more
  \item \emph{Say it:} \emph{wǒmen zuò kèchē qù nàr} — We're going there by coach.
  \item \emph{Recall:} say the Mandarin for public transport, then the Mandarin for to drop coins in the box, then say \emph{kèchē} again
\end{itemize}

\begin{tcolorbox}[breakable,colback=teal!4,colframe=teal!35!black,title={Before you move on}]
What does \zh{客车} mean? (\textquotedblleft{}Coach; long-distance bus\textquotedblright{}.) Say it once more.
\end{tcolorbox}
\section[\texorpdfstring{\zh{起点} (qǐdiǎn)}{起点 (qǐdiǎn)}]{\zh{起点} (qǐdiǎn) — starting point}

\label{lesson:ZH-C408-qidian}



\begin{quote}
Before the new one: say the Mandarin for to drop coins in the box, then the Mandarin for a long-distance coach.
\end{quote}

\subsection*{You'll want to know: \zh{起点}}
\textbf{\zh{起点}} — \emph{qǐdiǎn} — \textquotedblleft{}starting point\textquotedblright{}.

\begin{quote}
\textbf{\zh{公交的起点在火车站}} \emph{gōngjiāo de qǐdiǎn zài huǒchēzhàn} — The bus route starts at the railway station.
\end{quote}

Two third tones meet in this word: qǐ is said with a rising tone before diǎn, and the spelling keeps both marks.

\begin{cousinweb}
\textbf{\zh{起点}} (\emph{qǐdiǎn}) is built from \zh{起} 'to rise, to start' and \zh{点} 'a point', both inherited from Middle Chinese, \emph{khiX} and \emph{temX} (Guangyun, in Baxter's transcription), which glosses the first 'to rise' and the second as a dot.
\end{cousinweb}

\begin{grammarlens}[title={What you've built}]
One more word for getting around.
\end{grammarlens}

\subsection*{Your turn}
\begin{itemize}
\raggedright
  \item \emph{Say it:} \emph{qǐdiǎn}
  \item \emph{Say it:} \emph{qǐdiǎn}, once more
  \item \emph{Say it:} \emph{gōngjiāo de qǐdiǎn zài huǒchēzhàn} — The bus route starts at the railway station.
  \item \emph{Recall:} say the Mandarin for to drop coins in the box, then the Mandarin for a long-distance coach, then say \emph{qǐdiǎn} again
\end{itemize}

\begin{tcolorbox}[breakable,colback=teal!4,colframe=teal!35!black,title={Before you move on}]
What does \zh{起点} mean? (\textquotedblleft{}Starting point\textquotedblright{}.) Say it once more.
\end{tcolorbox}
